\section{Tema d'esame del 12 Gennaio 2026}
\subsection*{Esercizio 1}


Si traducano in logica del prim'ordine con uguaglianza le seguenti affermazioni in linguaggio naturale, utilizzando solo i simboli costanti presenti nelle frasi e i simboli di funzione e di predicato indicati di seguito. Simboli di funzione: \(s(\cdot) = \) \q{il successore di}; \(exp(2,\cdot) = \) \q{la potenza \(\cdot\)-esima di 2}. Simboli di predicato: \(N(\cdot) = \) \q{\(\cdot\) è un numero}; \(P(\cdot) = \) \q{\(\cdot\) è un numero pari}; \(D(\cdot) = \) \q{\(\cdot\) è un numero dispari}; \(Pr(\cdot) = \) \q{\(\cdot\) è un numero primo}; \(Div(\cdot,\cdot) = \) \q{\(\cdot\) è divisibile per \(\cdot\)}; \((\cdot \le \cdot) = \) \q{minore o uguale a}.
\paragraph{Svolgimento:}
\begin{itemize}
    \item \textbf{(a)} I numeri primi sono divisibili solo per 1 e per loro stessi.
    \[\forall x ((N(x) \land Pr(x)) \to \forall y ((N(y) \land Div(x,y)) \to (y=1 \lor y=x)))\]

    \item \textbf{(b)} I successori dei numeri dispari sono tutti divisibili per lo stesso numero.
    \[\exists d \st (N(d) \land \forall x ((N(x) \land D(x)) \to Div(s(x), d)))\]

    \item \textbf{(c)} Ciascun numero primo è compreso tra due potenze di 2 adiacenti.
    \[\forall x ((N(x) \land Pr(x)) \to \exists i \st (N(i) \land (exp(2,i) \le x \land x \le exp(2,s(i)))))\]

    \item \textbf{(d)} Qualche numero pari è divisibile da non più di 3 numeri.
    \[
    \begin{aligned}
      \exists x& \st (N(x) \land P(x) \land \forall i \forall j \forall k \forall l ((N(i) \land N(j) \land \\
      & N(k) \land N(l) \land Div(x,i) \land Div(x,j) \land Div(x,k) \land Div(x,l)) \to\\
      & (i=j \lor i=k \lor i=l \lor j=k \lor j=l \lor k=l)))
    \end{aligned}
    \]

    \item \textbf{(e)} Ci sono infiniti numeri primi.
    \[\forall x ((N(x) \land Pr(x)) \to \exists y \st (N(y) \land Pr(y) \land \neg(y \le x)))\]
\end{itemize}


\subsection*{Esercizio 2}


Si indichi, per ciascuna delle seguenti formule, se sono valide, insoddisfacibili o solo soddisfacibili. In ciascun caso, motivare semanticamente la risposta data. Esibire, successivamente, una deduzione in Natural Deduction per ogni formula che risulta valida o per la sua negazione, se la formula risulta insoddisfacibile.


\subsubsection*{(a) \(((\neg Q \lor P) \lor (\neg R \land P)) \to (Q \to P)\)}
\paragraph{Svolgimento:}
\textbf{Valida.} Semanticamente, la premessa si semplifica in \((\neg Q \lor P) \lor (\neg R \land P)\). Poiché \(\neg Q \lor P \equiv Q \to P\), l'intera implicazione ha la forma \((A \lor B) \to A\), che se analizzata logicamente rende sempre vero il conseguente se la disgiunzione è vera, nel caso si derivi dalla prima parte, altrimenti da \(\neg R \land P\) si deriva \(P\), da cui segue banalmente \(Q \to P\)[cite: 1].

\begin{prooftree}
  \small
  \AxiomC{\({[((\neg Q \lor P) \lor (\neg R \land P))]}_1\)}
  
  \AxiomC{\({[\neg Q \lor P]}_2\)}
  \AxiomC{\({[\neg Q]}_4\)}
  \AxiomC{\({[Q]}_3\)}
  \RuleLabel{\(\neg E\)}
  \BinaryInfC{\(\bot\)}
  \RuleLabel{\(\bot E\)}
  \UnaryInfC{\(P\)}
  \AxiomC{\({[P]}_4\)}
  \RuleLabel{\(\lor E_{4}\)}
  \TrinaryInfC{\(P\)}
    
  \AxiomC{\({[\neg R \land P]}_2\)}
  \RuleLabel{\(\land E\)}
  \UnaryInfC{\(P\)}
  
  \RuleLabel{\(\lor E_{2}\)}
  \TrinaryInfC{\(P\)}
  \RuleLabel{\(\to I_3\)}
  \UnaryInfC{\(Q \to P\)}
  \RuleLabel{\(\to I_1\)}
  \UnaryInfC{\(((\neg Q \lor P) \lor (\neg R \land P)) \to (Q \to P)\)}
\end{prooftree}

\subsubsection*{(b) \(\neg\neg(P \lor \neg P)\)}
\paragraph{Svolgimento:}
\textbf{Valida.} Equivalenza del terzo escluso[cite: 1].
\begin{prooftree}
  \AxiomC{\({[\neg (P \lor \neg P)]}_1\)}
    \AxiomC{\({[P]}_2\)}
    \RuleLabel{\(\lor I\)}
    \UnaryInfC{\(P \lor \neg P\)}
    \RuleLabel{\(\neg E\)}
    \BinaryInfC{\(\bot\)}
    \RuleLabel{\(\neg I_2\)}
    \UnaryInfC{\(\neg P\)}
    \RuleLabel{\(\lor I\)}
    \UnaryInfC{\(P \lor \neg P\)}
    \AxiomC{\({[\neg (P \lor \neg P)]}_1\)}
    \RuleLabel{\(\neg E\)}
    \BinaryInfC{\(\bot\)}
    \RuleLabel{\(\neg I_1\)}
    \UnaryInfC{\(\neg\neg(P \lor \neg P)\)}
\end{prooftree}

\subsubsection*{(c) \((P \lor Q) \to (\neg P \to Q)\)}
\paragraph{Svolgimento:}
\textbf{Valida.} Riscrittura diretta del sillogismo disgiuntivo[cite: 1].
\begin{prooftree}
    \AxiomC{\({[P \lor Q]}_1\)}
    \AxiomC{\({[P]}_2\)}
    \AxiomC{\({[\neg P]}_3\)}
    \RuleLabel{\(\neg E\)}
    \BinaryInfC{\(\bot\)}
    \RuleLabel{\(\bot E\)}
    \UnaryInfC{\(Q\)}
    \AxiomC{\({[Q]}_2\)}
    \RuleLabel{\(\lor E_{2}\)}
    \TrinaryInfC{\(Q\)}
    \RuleLabel{\(\to I_3\)}
    \UnaryInfC{\(\neg P \to Q\)}
    \RuleLabel{\(\to I_1\)}
    \UnaryInfC{\((P \lor Q) \to (\neg P \to Q)\)}
\end{prooftree}

\subsubsection*{(d) \((P \to (Q \lor \neg R)) \lor (( \neg P \land \neg Q ) \to \neg R)\)}
\paragraph{Svolgimento:}
\textbf{Valida.} Si dimostra per assurdo, con una costruzione strutturata che ricava dapprima \(\neg P\) tramite \(\RAA\) e successivamente costruisce l'implicazione di sinistra.
\paragraph{Sottoalbero \(1\)}
\begin{prooftree}
    \tiny
    \AxiomC{\({[\neg ((P \to (Q \lor \neg R)) \lor ((\neg P \land \neg Q) \to \neg R))]}_1\)}
    
    \AxiomC{\({[P]}_2\)}
    \AxiomC{\({[\neg P \land \neg Q]}_3\)}
    \RuleLabel{\(\land E\)}
    \UnaryInfC{\(\neg P\)}
    \RuleLabel{\(\neg E\)}
    \BinaryInfC{\(\bot\)}
    \RuleLabel{\(\bot E\)}
    \UnaryInfC{\(\neg R\)}
    \RuleLabel{\(\to I_3\)}
    \UnaryInfC{\((\neg P \land \neg Q) \to \neg R\)}
    \RuleLabel{\(\lor I\)}
    \UnaryInfC{\((P \to (Q \lor \neg R)) \lor ((\neg P \land \neg Q) \to \neg R)\)}
    
    \RuleLabel{\(\neg E\)}
    \BinaryInfC{\(\bot\)}
    \RuleLabel{\(\RAA_2\)}
    \UnaryInfC{\(\neg P\)}
    
    \AxiomC{\({[P]}_4\)}
    \RuleLabel{\(\neg E\)}
    \BinaryInfC{\(\bot\)}
    \RuleLabel{\(\bot E\)}
    \UnaryInfC{\(Q \lor \neg R\)}
\end{prooftree}
\begin{prooftree}
  \AxiomC{Sottoalbero \(1\)}
    \UnaryInfC{\(Q \lor \neg R\)}
    \RuleLabel{\(\to I_4\)}
    \UnaryInfC{\(P \to (Q \lor \neg R)\)}
    \RuleLabel{\(\lor I\)}
    \UnaryInfC{\((P \to (Q \lor \neg R)) \lor ((\neg P \land \neg Q) \to \neg R)\)}
    
    \AxiomC{\({[\neg ((P \to (Q \lor \neg R)) \lor ((\neg P \land \neg Q) \to \neg R))]}_1\)}
    \RuleLabel{\(\neg E\)}
    \BinaryInfC{\(\bot\)}
    
    \RuleLabel{\(\RAA_1\)}
    \UnaryInfC{\((P \to (Q \lor \neg R)) \lor ((\neg P \land \neg Q) \to \neg R)\)}
\end{prooftree}

\subsubsection*{(e) \((((\neg Q \lor \neg R) \to \neg P) \lor (P \to (Q \lor R))) \to (R \to P)\)}
\paragraph{Svolgimento:}
\textbf{Solo soddisfacibile.} Se poniamo \(R = \top\) e \(P = \bot\), la premessa globale valuta a vero e la conclusione \(R \to P\) valuta a falso, rendendo l'implicazione complessivamente falsa. Se invece poniamo \(P = \top\), la formula valuta a vero[cite: 1]. Non essendoci un risultato di validità o insoddisfacibilità assoluta, non presentiamo una deduzione.


\subsection*{Esercizio 3}


Si indichi, per ciascuna delle seguenti formule se sono valide, insoddisfacibili o solo soddisfacibili. Motivare semanticamente. Esibire una deduzione in Natural Deduction per le valide o per la negazione delle insoddisfacibili.


\paragraph{Svolgimento:}

\subsubsection*{(a) \(\neg\exists x \st (A(x) \lor B(x)) \to \neg\neg\forall x \st \neg(A(x) \lor B(x))\)}
\paragraph{Svolgimento:}
\textbf{Valida.}
\begin{prooftree}
    \AxiomC{\({[\neg \exists x \st (A(x) \lor B(x))]}_1\)}
    \AxiomC{\({[A(y) \lor B(y)]}_2\)}
    \RuleLabel{\(\exists I\)}
    \UnaryInfC{\(\exists x \st (A(x) \lor B(x))\)}
    \RuleLabel{\(\neg E\)}
    \BinaryInfC{\(\bot\)}
    \RuleLabel{\(\neg I_2\)}
    \UnaryInfC{\(\neg (A(y) \lor B(y))\)}
    \RuleLabel{\(\forall I\)}
    \UnaryInfC{\(\forall x \st \neg (A(x) \lor B(x))\)}
    
    \AxiomC{\({[\neg \forall x \st \neg (A(x) \lor B(x))]}_3\)}
    \RuleLabel{\(\neg E\)}
    \BinaryInfC{\(\bot\)}
    \RuleLabel{\(\neg I_3\)}
    \UnaryInfC{\(\neg\neg\forall x \st \neg (A(x) \lor B(x))\)}
    
    \RuleLabel{\(\to I_1\)}
    \UnaryInfC{\(\neg \exists x \st (A(x) \lor B(x)) \to \neg\neg\forall x \st \neg(A(x) \lor B(x))\)}
\end{prooftree}

\subsubsection*{(b) \((\exists x \st A(x) \to \forall x \st B(x)) \to \forall x \st (A(x) \to B(x))\)}

\paragraph{Svolgimento:}
\textbf{Valida.}
\begin{prooftree}
    \AxiomC{\({[\exists x \st A(x) \to \forall x \st B(x)]}_1\)}
    \AxiomC{\({[A(y)]}_2\)}
    \RuleLabel{\(\exists I\)}
    \UnaryInfC{\(\exists x \st A(x)\)}
    \RuleLabel{\(\to E\)}
    \BinaryInfC{\(\forall x \st B(x)\)}
    \RuleLabel{\(\forall E\)}
    \UnaryInfC{\(B(y)\)}
    \RuleLabel{\(\to I_2\)}
    \UnaryInfC{\(A(y) \to B(y)\)}
    \RuleLabel{\(\forall I\)}
    \UnaryInfC{\(\forall x \st (A(x) \to B(x))\)}
    
    \RuleLabel{\(\to I_1\)}
    \UnaryInfC{\((\exists x \st A(x) \to \forall x \st B(x)) \to \forall x \st (A(x) \to B(x))\)}
\end{prooftree}

\subsubsection*{(c) \(((b=a) \land (\exists x \st C(a,x)) \land \forall x \forall y (C(b,x) \to C(b,y))) \to \exists x \st C(b,x)\)}
\paragraph{Svolgimento:}
\textbf{Valida.}
\begin{prooftree}
    \tiny
    \AxiomC{\({[(b=a) \land (\exists x \st C(a,x)) \land \forall x \forall y (C(b,x) \to C(b,y))]}_1\)}
    \RuleLabel{\(\land E\)}
    \UnaryInfC{\(b = a\)}
    \RuleLabel{\(Symm\)}
    \UnaryInfC{\(a = b\)}
    \RuleLabel{\(Sub_{\exists x \st C(z,x)}\)}
    \UnaryInfC{\(\exists x \st C(a,x) \to \exists x \st C(b,x)\)}
    
    \AxiomC{\({[(b=a) \land (\exists x \st C(a,x)) \land \forall x \forall y (C(b,x) \to C(b,y))]}_1\)}
    \RuleLabel{\(\land E\)}
    \UnaryInfC{\(\exists x \st C(a,x)\)}
    
    \RuleLabel{\(\to E\)}
    \BinaryInfC{\(\exists x \st C(b,x)\)}
    
    \RuleLabel{\(\to I_1\)}
    \UnaryInfC{\(((b=a) \land (\exists x \st C(a,x)) \land \forall x \forall y (C(b,x) \to C(b,y))) \to \exists x \st C(b,x)\)}
\end{prooftree}

\subsubsection*{(d) \((\forall x \st A(x) \to \exists x \st B(x)) \to \exists x \st (A(x) \to B(x))\)}
\paragraph{Svolgimento:}
\textbf{Valida.}
\begin{prooftree}
    \tiny
    \AxiomC{\({[\forall x \st A(x) \to \exists x \st B(x)]}_1\)}
    
    \AxiomC{\({[\neg \exists x \st (A(x) \to B(x))]}_2\)}
    \AxiomC{\({[\neg A(y)]}_3\)}
    \AxiomC{\({[A(y)]}_4\)} 
    \RuleLabel{\(\neg E\)} 
    \BinaryInfC{\(\bot\)}
    \RuleLabel{\(\bot E\)} 
    \UnaryInfC{\(B(y)\)}
    \RuleLabel{\(\to I_4\)}
    \UnaryInfC{\(A(y) \to B(y)\)}
    \RuleLabel{\(\exists I\)}
    \UnaryInfC{\(\exists x \st (A(x) \to B(x))\)}
    \RuleLabel{\(\neg E\)}
    \BinaryInfC{\(\bot\)}
    \RuleLabel{\(\RAA_3\)}
    \UnaryInfC{\(A(y)\)}
    \RuleLabel{\(\forall I\)}
    \UnaryInfC{\(\forall x \st A(x)\)}
    
    \RuleLabel{\(\to E\)}
    \BinaryInfC{\(\exists x \st B(x)\)}
    
    \AxiomC{\({[\neg \exists x \st (A(x) \to B(x))]}_2\)}
    \AxiomC{\({[B(z)]}_5\)}
    \RuleLabel{\(\to I_6\)}
    \UnaryInfC{\(A(z) \to B(z)\)}
    \RuleLabel{\(\exists I\)}
    \UnaryInfC{\(\exists x \st (A(x) \to B(x))\)}
    \RuleLabel{\(\neg E\)}
    \BinaryInfC{\(\bot\)}
    
    \RuleLabel{\(\exists E_5\)}
    \BinaryInfC{\(\bot\)}
    
    \RuleLabel{\(\RAA_2\)}
    \UnaryInfC{\(\exists x \st (A(x) \to B(x))\)}
    
    \RuleLabel{\(\to I_1\)}
    \UnaryInfC{\((\forall x \st A(x) \to \exists x \st B(x)) \to \exists x \st (A(x) \to B(x))\)}
\end{prooftree}

\subsubsection*{(e) \(\exists x \st \forall y ((\neg C(x,x) \to C(x,y)) \land (C(x,y) \to \neg C(x,x)))\)}
\paragraph{Svolgimento:}
\textbf{Insoddisfacibile.} Semanticamente, ponendo \(y=x\), otteniamo l'equivalenza \(C(x,x) \leftrightarrow \neg C(x,x)\), che è una palese contraddizione per ogni modello. Si mostra la dimostrazione della sua negazione.
\paragraph{Sottoalbero 1}
\begin{prooftree}
    \tiny
    
    \AxiomC{\({[\forall y ((\neg C(z,z) \to C(z,y)) \land (C(z,y) \to \neg C(z,z)))]}_2\)}
    \RuleLabel{\(\forall E\)}
    \UnaryInfC{\((\neg C(z,z) \to C(z,z)) \land (C(z,z) \to \neg C(z,z))\)}
    \RuleLabel{\(\land E\)}
    \UnaryInfC{\(C(z,z) \to \neg C(z,z)\)}
    \AxiomC{\({[C(z,z)]}_3\)}
    \RuleLabel{\(\to E\)}
    \BinaryInfC{\(\neg C(z,z)\)}
    \AxiomC{\({[C(z,z)]}_3\)}
    \RuleLabel{\(\neg E\)}
    \BinaryInfC{\(\bot\)}
    \RuleLabel{\(\neg I_3\)}
    \UnaryInfC{\(\neg C(z,z)\)}
    \end{prooftree}
    \paragraph{Sottoalbero 2}
    \begin{prooftree}
      \tiny
      
      \AxiomC{\({[\forall y ((\neg C(z,z) \to C(z,y)) \land (C(z,y) \to \neg C(z,z)))]}_2\)}
      \RuleLabel{\(\forall E\)}
    \UnaryInfC{\((\neg C(z,z) \to C(z,z)) \land (C(z,z) \to \neg C(z,z))\)}
    \RuleLabel{\(\land E\)}
    \UnaryInfC{\(\neg C(z,z) \to C(z,z)\)}
    
    \AxiomC{\({[\forall y ((\neg C(z,z) \to C(z,y)) \land (C(z,y) \to \neg C(z,z)))]}_2\)}
    \RuleLabel{\(\forall E\)}
    \UnaryInfC{\((\neg C(z,z) \to C(z,z)) \land (C(z,z) \to \neg C(z,z))\)}
    \RuleLabel{\(\land E\)}
    \UnaryInfC{\(C(z,z) \to \neg C(z,z)\)}
    \AxiomC{\({[C(z,z)]}_4\)}
    \RuleLabel{\(\to E\)}
    \BinaryInfC{\(\neg C(z,z)\)}
    \AxiomC{\({[C(z,z)]}_4\)}
    \RuleLabel{\(\neg E\)}
    \BinaryInfC{\(\bot\)}
    \RuleLabel{\(\neg I_4\)}
    \UnaryInfC{\(\neg C(z,z)\)}
    
    \RuleLabel{\(\to E\)}
    \BinaryInfC{\(C(z,z)\)}
  \end{prooftree}
  \begin{prooftree}
    \AxiomC{Sottoalbero 1}
    \AxiomC{Sottoalbero 2}

    \RuleLabel{\(\neg E\)}
    \BinaryInfC{\(\bot\)}
    
    \AxiomC{\({[\exists x \st \forall y ((\neg C(x,x) \to C(x,y)) \land (C(x,y) \to \neg C(x,x)))]}_1\)}
    \RuleLabel{\(\exists E_2\)}
    \BinaryInfC{\(\bot\)}
    \RuleLabel{\(\neg I_1\)}
    \UnaryInfC{\(\neg \exists x \st \forall y ((\neg C(x,x) \to C(x,y)) \land (C(x,y) \to \neg C(x,x)))\)}
\end{prooftree}